\documentclass[12pt,a4paper]{article}
\usepackage[T2A]{fontenc}
\usepackage[utf8]{inputenc}
\usepackage[russian]{babel}
\usepackage{geometry}
\geometry{top=2cm,bottom=2cm,left=2.5cm,right=1.5cm}
\usepackage{amsmath}
\usepackage{indentfirst}
\setlength{\parindent}{1.25cm}

\begin{document}

% ========== ТИТУЛЬНЫЙ ЛИСТ ==========
\begin{titlepage}
    \begin{center}
        \textbf{БЮДЖЕТНОЕ УЧРЕЖДЕНИЕ ВЫСШЕГО ОБРАЗОВАНИЯ}\\
        \textbf{ХАНТЫ-МАНСИЙСКОГО АВТОНОМНОГО ОКРУГА – ЮГРЫ}\\
        \textbf{«СУРГУТСКИЙ ГОСУДАРСТВЕННЫЙ УНИВЕРСИТЕТ»}
        
        \vspace{0.8cm}
        ПОЛИТЕХНИЧЕСКИЙ ИНСТИТУТ
        
        \vspace{0.2cm}
        Кафедра прикладной математики
        
        \vspace{3.5cm}
        \textbf{КУРСОВАЯ РАБОТА}
        
        \vspace{0.3cm}
        по дисциплине «Численные методы»
        
        \vspace{2.5cm}
        ТЕМА КУРСОВОЙ РАБОТЫ
        
        \vspace{0.5cm}
        \underline{\hspace{8cm}}\\
        \vspace{0.2cm}
        {\large \textbf{Метод конечных разностей решения линейной}} \\
        {\large \textbf{двухточечной краевой задачи для ОДУ второго порядка}}
    \end{center}
    
    \vspace{2.5cm}
    \begin{flushright}
        \begin{minipage}{0.6\textwidth}
            \raggedright
            Выполнил \\
            Студент гр. 601-\_\_1 \hspace{0.8cm} \underline{\hspace{2cm}} \hspace{0.5cm} / ФИО. \\
            \hfill (подпись)
            
            \vspace{0.8cm}
            Проверил \\
            \underline{\hspace{3cm}} \hspace{0.8cm} (оценка) \\
            К.ф.-м.н, доцент \hspace{0.8cm} \underline{\hspace{2cm}} \hspace{0.5cm} / Назин А.Г. \\
            \hfill (подпись)
        \end{minipage}
    \end{flushright}
    
    \vfill
    \begin{center}
        Сургут 2026
    \end{center}
\end{titlepage}

% ========== ВВЕДЕНИЕ ==========
\section*{Введение}
\addcontentsline{toc}{section}{Введение}

Многие физические, технические и экономические процессы описываются дифференциальными уравнениями. Особый класс составляют \textbf{двухточечные краевые задачи} для обыкновенных дифференциальных уравнений (ОДУ) (то есть уравнений, связывающих функцию и её производные, при этом функция зависит только от одной переменной) второго порядка. Такие задачи возникают, когда известны условия не в одной точке, а на двух концах некоторого отрезка — отсюда название «двухточечные».

Математическая постановка задачи, рассматриваемой в данной работе, имеет вид:

\begin{equation*}
    -\frac{d}{dx}\left(k(x)\frac{du}{dx}\right) + q(x)u = f(x), \quad a < x < b
\end{equation*}

Здесь $u(x)$ — искомая функция, $k(x)$, $q(x)$, $f(x)$ — заданные функции координаты, причём $k(x) > 0$.

На границах отрезка $[a, b]$ задаются дополнительные условия — первого, второго или третьего рода.

Аналитическое решение такой задачи удаётся найти лишь для простейших случаев. Для реальных задач с переменными коэффициентами $k(x)$, $q(x)$, $f(x)$ приходится использовать \textbf{численные методы}. Один из самых распространённых — \textbf{метод конечных разностей} (МКР). Суть метода заключается в замене непрерывной области $[a, b]$ набором отдельных точек (узлов), а производных — приближёнными формулами через значения функции в этих точках. В результате дифференциальное уравнение превращается в систему линейных алгебраических уравнений, которую можно решить на компьютере.

\textbf{Цель работы:}
\begin{enumerate}
    \item Описать постановку линейной двухточечной краевой задачи для ОДУ второго порядка с разными граничными условиями.
    \item Указать условия разрешимости краевой задачи.
    \item Описать метод конечных разностей для решения этой задачи.
    \item Реализовать в среде Octave алгоритм решения краевой задачи методом конечных разностей для одномерного стационарного уравнения теплопроводности
    \begin{equation*}
        -\frac{d}{dx}\left(k(x)\frac{du}{dx}\right) + q(x)u = f(x), \quad a < x < b,
    \end{equation*}
    где $k(x) \neq \text{const}$, $q(x) \neq \text{const}$, $f(x) \neq \text{const}$ — непрерывные функции, вид которых выбирается самостоятельно. Требуемая точность — $10^{-4}$.
    \item Оценить погрешность численного решения по правилу Рунге.
    \item Провести вычислительные эксперименты с различными краевыми условиями.
    \item Построить графики полученного решения и график погрешности.
\end{enumerate}

\newpage

% ========== ПОСТАНОВКА КРАЕВОЙ ЗАДАЧИ ==========
\section*{Постановка краевой задачи}
\addcontentsline{toc}{section}{Постановка краевой задачи}

Рассмотрим линейную двухточечную краевую задачу для обыкновенного дифференциального уравнения второго порядка на отрезке $[a, b]$:

\begin{equation}
    -\frac{d}{dx}\left(k(x)\frac{du}{dx}\right) + q(x)u = f(x), \quad a < x < b
\end{equation}

где $u(x)$ — искомая функция, $k(x)$, $q(x)$, $f(x)$ — заданные непрерывные функции, причём $k(x) > 0$ на всём отрезке $[a, b]$.

Для единственности решения к уравнению (1) необходимо добавить граничные условия на концах отрезка. В зависимости от физической постановки могут задаваться условия первого, второго или третьего рода.

\subsection*{Граничные условия первого рода (условия Дирихле)}

На границах отрезка задаются значения самой искомой функции:

\begin{equation}
    u(a) = \alpha, \qquad u(b) = \beta
\end{equation}

где $\alpha$, $\beta$ — заданные числа.

\subsection*{Граничные условия второго рода (условия Неймана)}

На границах отрезка задаются значения производной искомой функции:

\begin{equation}
    u'(a) = \gamma, \qquad u'(b) = \delta
\end{equation}

где $\gamma$, $\delta$ — заданные числа.

\subsection*{Граничные условия третьего рода (смешанные)}

На границах отрезка задаётся линейная комбинация функции и её производной:

\begin{equation}
    \alpha_0 u(a) + \alpha_1 u'(a) = \gamma_a, \qquad \beta_0 u(b) + \beta_1 u'(b) = \gamma_b
\end{equation}

где $\alpha_0, \alpha_1, \beta_0, \beta_1, \gamma_a, \gamma_b$ — заданные числа, причём $|\alpha_0| + |\alpha_1| \neq 0$ и $|\beta_0| + |\beta_1| \neq 0$.

\subsection*{Условия разрешимости}

Для существования и единственности решения краевой задачи (1) необходимо выполнение определённых условий.

\begin{itemize}
    \item Если $q(x) \geq 0$ на $[a, b]$ и хотя бы в одной точке $q(x) > 0$, то задача с граничными условиями первого или третьего рода имеет единственное решение.
    \item Если $q(x) \equiv 0$ на $[a, b]$, то решение существует и единственно при выполнении дополнительных условий согласования (например, для задачи с условиями Неймана необходимо, чтобы интеграл от правой части $f(x)$ был равен нулю).
    \item Условие $k(x) > 0$ гарантирует, что уравнение не вырождается.
\end{itemize}

\end{document}